\PassOptionsToPackage{table,xcdraw}{xcolor}
\documentclass[10pt, a4paper]{article}

\usepackage[T1]{fontenc}
\usepackage[utf8]{inputenc}
\usepackage{tgpagella}
\usepackage[spanish, es-noshorthands]{babel}
\usepackage{amsmath, amssymb}
\usepackage{booktabs}
\usepackage{tabularx}
\usepackage{xcolor}
\usepackage{tikz}
\usetikzlibrary{shapes.geometric, arrows.meta, positioning}
\usepackage[most]{tcolorbox}
\usepackage[margin=1.6cm]{geometry}
\usepackage{fancyhdr}

\setlength{\headheight}{16pt}
\setlength{\emergencystretch}{2em}
\definecolor{ucbblue}{RGB}{0, 51, 102}

\pagestyle{fancy}
\fancyhf{}
\lhead{\textbf{IMT-221}}
\chead{Checkpoint de Integración y Readiness: 21 de Octubre (Hito 2)}
\rhead{W12S01}
\lfoot{Rumbo a Defensa Hito 2}
\rfoot{Página \thepage}
\renewcommand{\headrulewidth}{0.4pt}

\begin{document}

\begin{center}
    \Large\textbf{\textcolor{ucbblue}{Hoja de Preparación y Registro: Checkpoint del 21 de Octubre (Hito 2)}}
\end{center}

\noindent\textbf{Grupo:} \rule{3cm}{0.4pt} \hfill \textbf{Integrantes:} \rule{7.5cm}{0.4pt} \hfill \textbf{Fecha Registro Inicial:} 19/10/2026

\section*{1. Comparación Sistémica Hito 1 vs. Hito 2 en Nodos Homólogos (Visual G)}
\textit{El Hito 2 exige demostrar una mejora de sistema bajo condiciones equivalentes, no simplemente mostrar que el amplificador tiene ganancia.}

\begin{center}
    \begin{tikzpicture}[>=Stealth, thick, scale=0.88, transform shape]
        % --- FILA SUPERIOR: HITO 1 (PASIVO) ---
        \node[anchor=east, font=\bfseries] at (-0.2, 0) {Baseline Hito 1 (Pasivo):};
        
        \node[draw, rectangle, rounded corners=2pt, fill=gray!15, minimum height=0.9cm, text width=2.5cm, align=center] (b1) at (1.4, 0) {Puente CA /\\Sensor NTC};
        
        \node[draw, rectangle, rounded corners=2pt, fill=yellow!20, minimum height=0.9cm, text width=2.7cm, align=center] (t1) at (6.4, 0) {Twin-T $50\,\text{Hz}$\\(Carga Pasiva)};
        
        % Conexion y Nodo Homologo N_A (con linea indicadora superior libre de solapamiento)
        \draw[->] (b1.east) -- (t1.west);
        \filldraw[red!80!black] (3.9, 0) circle (2.2pt) coordinate (na1);
        \draw[red!80!black, thin] (na1) -- (3.9, 0.65) node[above, font=\footnotesize, fill=red!5, draw=red!40, rounded corners=1pt, inner sep=2pt] {Nodo $N_A$: $V_{\mathrm{nodo,H1}}$};

        % Salida Hito 1
        \draw[->] (t1.east) -- (8.8, 0);
        \filldraw[blue!80!black] (8.8, 0) circle (2.2pt);
        \node[anchor=west, font=\footnotesize] at (9.0, 0) {\textbf{Salida H1} ($V_{u,H1},\; V_{50,H1}$)};

        % --- FILA INFERIOR: HITO 2 (ACTIVO) ---
        \node[anchor=east, font=\bfseries] at (-0.2, -2.1) {Sistema Hito 2 (Activo):};
        
        \node[draw, rectangle, rounded corners=2pt, fill=gray!15, minimum height=0.9cm, text width=2.5cm, align=center] (b2) at (1.4, -2.1) {Puente CA /\\Sensor NTC};
        
        \node[draw, rectangle, rounded corners=2pt, fill=ucbblue!15, minimum height=0.9cm, text width=2.7cm, align=center] (in2) at (6.4, -2.1) {In-Amp Activo\\(Alta $Z_{in}$)};
        
        \node[draw, rectangle, rounded corners=2pt, fill=green!15, minimum height=0.9cm, text width=2.4cm, align=center] (t2) at (9.9, -2.1) {Etapa Filtro\\Twin-T $50\,\text{Hz}$};

        % Conexion y Nodo Homologo N_A' (con linea indicadora inferior libre de solapamiento)
        \draw[->] (b2.east) -- (in2.west);
        \filldraw[red!80!black] (3.9, -2.1) circle (2.2pt) coordinate (na2);
        \draw[red!80!black, thin] (na2) -- (3.9, -2.75) node[below, font=\footnotesize, fill=red!5, draw=red!40, rounded corners=1pt, inner sep=2pt] {Nodo $N_A^\prime$: $V_{\mathrm{nodo,H2}}$};

        % Vinculo visual entre nodos homologos
        \draw[<->, dashed, red!70!black] (3.9, -0.2) -- node[fill=white, inner sep=1.5pt, font=\scriptsize\bfseries, text=red!80!black] {Mismo nodo ($K_{\mathrm{pres}}$)} (3.9, -1.9);

        % Conexion a filtro y Salida Hito 2
        \draw[->] (in2.east) -- (t2.west);
        \draw[->] (t2.east) -- (12.0, -2.1);
        \filldraw[blue!80!black] (12.0, -2.1) circle (2.2pt);
        \node[anchor=west, font=\footnotesize] at (12.2, -2.1) {\textbf{Salida H2} ($V_{u,H2},\; V_{50,H2}$)};
    \end{tikzpicture}
\end{center}

\subsection*{Opciones Matemáticas Condicionadas para Cuantificar la Mejora}
\begin{enumerate}
    \setlength{\itemsep}{1pt}
    \item \textbf{Métrica Original $K_L$ de Hito 1:} Si utiliza el factor $K_L$, recupere su definición exacta del informe de Hito 1 (no redefinir de memoria).
    \item \textbf{Cociente de Preservación en un Nodo Homólogo ($K_{\mathrm{pres}}$):} Midiendo el \textbf{mismo nodo del puente} ($N_A$ vs. $N_A^\prime$) bajo idéntica excitación:
    \begin{equation}
        \boxed{K_{\mathrm{pres}} = \frac{|V_{\mathrm{nodo,con\ H2}}|}{|V_{\mathrm{nodo,baseline\ H1}}|}}
    \end{equation}
    \textit{Ejemplo Ilustrativo 5 (No es dato de grupo):} Si el puente entregaba $70\,\text{mV}$ al conectarse a la red pasiva de Hito 1 y sube a $95\,\text{mV}$ al conectarse a la entrada de alta impedancia del In-Amp: $K_{\mathrm{pres}} = \frac{95}{70} \approx 1.36$.
    \item \textbf{Cambio Relativo de Amplitud Útil ($\eta_V$) y Relación Señal/Interferencia ($SIR$):}
    \begin{equation}
        \eta_V = \frac{|V_{u,H2}| - |V_{u,H1}|}{|V_{u,H1}|}\times 100\% \quad (V_{u,H1}\neq 0), \qquad SIR = 20\log_{10}\left(\frac{|V_u|}{|V_{50}|}\right), \qquad \boxed{\Delta SIR = SIR_{H2} - SIR_{H1}}
    \end{equation}
    \textit{Advertencia:} Un aumento en $\eta_V$ por sí solo refleja ganancia, \textbf{no demuestra menor efecto de carga ni mejor rechazo de interferencia}.
\end{enumerate}

\section*{2. Diseño del Experimento Comparativo (Completar el 19/oct para ejecutar el 21/oct)}
\begin{table}[htbp]
    \centering
    \footnotesize
    \renewcommand{\arraystretch}{1.3}
    \begin{tabularx}{\textwidth}{|>{\raggedright\arraybackslash}p{6.2cm}|>{\raggedright\arraybackslash}X|}
    \hline
    \rowcolor[HTML]{EFEFEF} 
    \textbf{Pregunta de Diseño Experimental} & \textbf{Definición Técnica del Grupo} \\ \hline
    1. Limitación documentada en Hito 1 & \\ \hline
    2. Métrica elegida y su ecuación exacta & \\ \hline
    3. Nodo y magnitud medida en el baseline (Hito 1) & \\ \hline
    4. Nodo homólogo que se medirá en Hito 2 & \\ \hline
    5. Condiciones que se mantendrán estrictamente iguales & \\ \hline
    6. Qué resultado experimental refutaría la mejora esperada & \\ \hline
    \end{tabularx}
\end{table}

\section*{3. Secuencia de Validación y Registro Prospectivo del Checkpoint 21/oct (Visual H)}
\begin{center}
    \begin{tikzpicture}[>=Stealth, thick, scale=0.85, transform shape]
        \node[draw, rectangle, rounded corners=2pt, fill=ucbblue!10, minimum height=0.95cm, text width=2.2cm, align=center, font=\footnotesize] (c1) at (0,0) {\textbf{1. Etapa Aislada}\\Validada};
        \node[draw, rectangle, rounded corners=2pt, fill=ucbblue!10, minimum height=0.95cm, text width=2.2cm, align=center, font=\footnotesize] (c2) at (3.1,0) {\textbf{2. Integración}\\Hito 1 + Hito 2};
        \node[draw, rectangle, rounded corners=2pt, fill=ucbblue!10, minimum height=0.95cm, text width=2.2cm, align=center, font=\footnotesize] (c3) at (6.2,0) {\textbf{3. Prueba de}\\Loading ($K_{\mathrm{pres}}$)};
        \node[draw, rectangle, rounded corners=2pt, fill=ucbblue!10, minimum height=0.95cm, text width=2.2cm, align=center, font=\footnotesize] (c4) at (9.3,0) {\textbf{4. Modo Común}\\y CMRR};
        \node[draw, rectangle, rounded corners=2pt, fill=ucbblue!10, minimum height=0.95cm, text width=2.2cm, align=center, font=\footnotesize] (c5) at (12.4,0) {\textbf{5. Prueba de}\\$50\,\text{Hz}$ / SIR};
        \node[draw, rectangle, rounded corners=2pt, fill=green!15, minimum height=0.95cm, text width=2.3cm, align=center, font=\footnotesize] (c6) at (15.6,0) {\textbf{6. Defensa 28/oct}\\(Condicionada)};
        
        \draw[->] (c1.east) -- (c2.west);
        \draw[->] (c2.east) -- (c3.west);
        \draw[->] (c3.east) -- (c4.west);
        \draw[->] (c4.east) -- (c5.west);
        \draw[->] (c5.east) -- (c6.west);
    \end{tikzpicture}
\end{center}

\begin{table}[htbp]
    \centering
    \footnotesize
    \renewcommand{\arraystretch}{1.28}
    \begin{tabularx}{\textwidth}{|>{\raggedright\arraybackslash}p{5.8cm}|>{\raggedright\arraybackslash}p{4.5cm}|>{\raggedright\arraybackslash}X|}
    \hline
    \rowcolor[HTML]{EFEFEF} 
    \textbf{Condición de Readiness (Verificar el 21/oct)} & \textbf{Estado / Ref. a Evidencia} & \textbf{Bloqueo Exacto y Próxima Acción} \\ \hline
    1. Gate A y diseño aprobados & & \\ \hline
    2. Etapa activa opera aislada con rieles válidos & & \\ \hline
    3. Ganancia diferencial medida y ensayo común/límite & & \\ \hline
    4. Baseline Hito 1 disponible e integración realizada & & \\ \hline
    5. Comparación de loading/métrica y prueba $50\,\text{Hz}$ & & \\ \hline
    6. Grounding revisado y discrepancia interpretada & & \\ \hline
    \end{tabularx}
\end{table}

\noindent\textbf{Clasificación del 21/oct (No marcar el 19/oct):} $\square$ \texttt{READY FOR DEFENSE} \quad $\square$ \texttt{MINOR EVIDENCE PENDING} \quad $\square$ \texttt{NOT READY} \\
\textbf{¿Se sostiene la defensa propuesta del 28/oct?} $\square$ Por confirmar \quad $\square$ Viable condicionado \quad $\square$ No viable por bloqueo técnico: \hrulefill

\end{document}
